% !TeX root = ../main.tex
\section*{บทคัดย่อ}
การจัดการลาของนิสิตอย่างเป็นระบบช่วยให้การยื่นคำขอ การพิจารณา
และการติดตามข้อมูลเป็นไปอย่างต่อเนื่องและตรวจสอบได้
ระบบลานิสิตสำหรับคณะวิศวกรรมศาสตร์
กำแพงแสน มหาวิทยาลัยเกษตรศาสตร์ จึงได้รับการพัฒนาขึ้นในรูปแบบเว็บแอปพลิเคชัน
เพื่อรองรับการยื่นคำขอและแนบหลักฐาน การพิจารณาตามลำดับการอนุมัติลา
และการติดตามสถานะ
โดยกำหนดสิทธิ์การเข้าถึงข้อมูลตามบทบาทและความสัมพันธ์กับนิสิต
การออกแบบระบบเชื่อมโยงข้อมูลรายวิชา ผู้สอน
และอาจารย์ที่ปรึกษาจากระบบข้อมูลสนับสนุน
พร้อมแยกคลังข้อมูลรายงานออกจากฐานข้อมูลธุรกรรม
เพื่อสรุปรายละเอียดต่างๆ สำหรับการลาให้เหมาะสมกับบทบาทของผู้ใช้
และรายวิชาที่ได้รับผลกระทบตามขอบเขตสิทธิ์ของผู้ใช้
\par\noindent\textbf{คำสำคัญ:} แอปพลิเคชัน, คลังข้อมูล, ลำดับการอนุมัติลา
\section*{Abstract}
Systematic student leave management supports continuity and traceability in
request submission, review, and tracking. A student leave management system
was therefore developed as a web application for the Faculty of Engineering
at Kamphaeng Saen, Kasetsart University. The system supports request
submission with supporting evidence, review according to the leave approval
sequence, and status tracking. Data access is governed by user roles and
relationships with students. The system design integrates subject, lecturer,
and advisor data from a supporting data management system and separates the
reporting warehouse from the operational database to summarize leave details
appropriate to each user's role and affected subjects within the user's
access scope.
\par\noindent\textbf{Keywords:} application, data warehouse, leave
approval workflow
